\begingroup\sloppy
\begin{references}
\item Drouin-Chartier JP, Chen S, Li Y, et al. Egg consumption and risk of cardiovascular disease: three large prospective US cohort studies, systematic review, and updated meta-analysis. BMJ. 2020;368:m513.
\item Centers for Disease Control NCHS. National Health and Nutrition Examination Survey (NHANES): History 2020 [cited 2020 Dec 15]. \url{https://www.cdc.gov/nchs/nhanes/history.htm}.
\item Andersen RE, Crespo CJ, Bartlett SJ, Cheskin LJ, Pratt M. Relationship of physical activity and television watching with body weight and level of fatness among children: results from the Third National Health and Nutrition Examination Survey. JAMA. 1998;279(12):938-942.
\item Sargent JD, Beach ML, Adachi-Mejia AM, et al. Exposure to movie smoking: its relation to smoking initiation among US adolescents. Pediatrics. 2005;116(5):1183-1191.
\item Morgenstern M, Sargent JD, Engels R, et al. Smoking in movies and adolescent smoking initiation: longitudinal study in six European countries. Am J Prev Med. 2013;44(4):339-344.
\item Dal Cin S, Stoolmiller M, Sargent JD. Exposure to smoking in movies and smoking initiation among black youth. Am J Prev Med. 2013;44(4):345-350.
\item Mejia R, Perez A, Pena L, et al. Smoking in movies and adolescent smoking initiation: a longitudinal study among Argentinian adolescents. J Pediatr. 2017;180:222-228.
\item Blackwell T, Yaffe K, Ancoli-Israel S, et al. Association of sleep characteristics and cognition in older community-dwelling men: the MrOS sleep study. Sleep. 2011;34(10):1347-1356.
\item Kalantar-Zadeh K, Abbott KC, Salahudeen AK, Kilpatrick RD, Horwich TB. Survival advantages of obesity in dialysis patients. Am J Clin Nutr. 2005;81(3):543-554.
\item McMillen RC, Gottlieb MA, Shaefer RM, Winickoff JP, Klein JD. Trends in electronic cigarette use among U.S. adults: use is increasing in both smokers and nonsmokers. Nicotine Tob Res. 2015;17(10):1195-1202.
\item Vandenbroucke JP. Prospective or retrospective: what’s in a name? BMJ. 1991;302(6771):249-250.
\item Rothman KJ, Greenland S, Lash TL. Modern Epidemiology. 3rd ed. Wolters Kluwer Health/Lippincott Williams \& Wilkins; 2008:758.
\sourcepage{239}
\item Bao Y, Bertoia ML, Lenart EB, et al. Origin, methods, and evolution of the three nurses’ health studies. Am J Public Health. 2016;106(9):1573-1581.
\item Man S, Xian Y, Holmes DN, et al. Association between thrombolytic door-to-needle time and 1-year mortality and readmission in patients with acute ischemic stroke. JAMA. 2020;323(21):2170-2184.
\item Browner WS, Pressman AR, Nevitt MC, Cauley JA, Cummings SR. Association between low bone density and stroke in elderly women: the study of osteoporotic fractures. Stroke. 1993;24(7):940-946.
\item Evans DS, Kim KM, Jiang X, Jacobson J, Browner W, Cummings SR. Prediction of in-hospital mortality among adults with COVID-19 infection. medRxiv. 2021.
\item Grodstein F, Stampfer MJ, Manson JE, et al. Postmenopausal estrogen and progestin use and the risk of cardiovascular disease. N Engl J Med. 1996;335(7):453-461.
\item Grodstein F, Manson JE, Colditz GA, Willett WC, Speizer FE, Stampfer MJ. A prospective, observational study of postmenopausal hormone therapy and primary prevention of cardiovascular disease. Ann Intern Med. 2000;133(12):933-941.
\item Hernán MA, Alonso A, Logan R, et al. Observational studies analyzed like randomized experiments: an application to postmenopausal hormone therapy and coronary heart disease. Epidemiology. 2008;19(6):766-779.
\item Zaroff JG, Cheetham TC, Palmetto N, et al. Association of azithromycin use with cardiovascular mortality. JAMA Netw Open. 2020;3(6):e208199.
\item Newman TB, Liljestrand P, Jeremy RJ, et al. Outcomes among newborns with total serum bilirubin levels of 25 mg per deciliter or more. N Engl J Med. 2006;354(18):1889-1900.
\item Escobar GJ, Liljestrand P, Hudes ES, et al. Five-year neurodevelopmental outcome of neonatal dehydration. J Pediatr. 2007;151(2):127-133, 133 e1.
\item Duarte D, El-Hagrassy MM, Couto TCE, Gurgel W, Fregni F, Correa H. Male and female physician suicidality: a systematic review and meta-analysis. JAMA Psychiatry. 2020;77(6):587-597.
\item Sterne JA, Hernán MA, Reeves BC, et al. ROBINS-I: a tool for assessing risk of bias in non-randomised studies of interventions. BMJ. 2016;355:i4919.
\item Hernán MA, Robins JM. Using big data to emulate a target trial when a randomized trial is not available. Am J Epidemiol. 2016;183(8):758-764.
\item Stefan MS, Shieh MS, Spitzer KA, et al. Association of antibiotic treatment with outcomes in patients hospitalized for an asthma exacerbation treated with systemic corticosteroids. JAMA Intern Med. 2019;179(3):333-339 [original version, subsequently retracted].
\item Newman TB. Possible immortal time bias in study of antibiotic treatment and outcomes in patients hospitalized for asthma. JAMA Intern Med. 2021;181(4):568-569.
\item Stefan MS, Shieh M-S, Spitzer KA, et al. Association of antibiotic treatment with outcomes in patients hospitalized for an asthma exacerbation treated with systemic corticosteroids. JAMA Intern Med. 2019;179(3):333-339 [replaced version of retracted article].
\item Newman T, Kohn M. Evidence-Based Diagnosis: An Introduction to Clinical Epidemiology. 2nd ed. Cambridge University Press; 2020:231-243.
\item Newman TB. Antibiotic treatment for inpatient asthma exacerbations: what have we learned? JAMA Intern Med. 2021;181(4):427-428.
\item Suissa S. Immortal time bias in pharmaco-epidemiology. Am J Epidemiol. 2008;167(4):492-499.
\item Hernán MA, Sauer BC, Hernandez-Diaz S, Platt R, Shrier I. Specifying a target trial prevents immortal time bias and other self-inflicted injuries in observational analyses. J Clin Epidemiol. 2016;79:70-75.
\item Vandenbroucke J, Pearce N. Point: incident exposures, prevalent exposures, and causal inference: does limiting studies to persons who are followed from first exposure onward damage epidemiology? Am J Epidemiol. 2015;182(10):826-833.
\item Hernán MA. Counterpoint: epidemiology to guide decision-making: moving away from practice-free research. Am J Epidemiol. 2015;182(10):834-839.
\item Vandenbroucke J, Pearce N. Vandenbroucke and Pearce respond to “incident and prevalent exposures and causal inference”. Am J Epidemiol. 2015;182(10):846-847.
\sourcepage{240}
\item Ray WA. Evaluating medication effects outside of clinical trials: new-user designs. Am J Epidemiol. 2003;158(9):915-920.
\item Simpson SH, Eurich DT, Majumdar SR, et al. A meta-analysis of the association between adherence to drug therapy and mortality. BMJ. 2006;333(7557):15.
\item Avins AL, Pressman A, Ackerson L, Rudd P, Neuhaus J, Vittinghoff E. Placebo adherence and its association with morbidity and mortality in the studies of left ventricular dysfunction. J Gen Intern Med. 2010;25(12):1275-1281.
\item Cummings SR, Lui LY, Eastell R, Allen IE. Association between drug treatments for patients with osteoporosis and overall mortality rates: a meta-analysis. JAMA Intern Med. 2019;179(11):1491-1500.
\item Hulley S, Grady D, Bush T, et al. Randomized trial of estrogen plus progestin for secondary prevention of coronary heart disease in postmenopausal women. Heart and Estrogen/progestin Replacement Study (HERS) Research Group. JAMA. 1998;280(7):605-613.
\item Bero L, Chartres N, Diong J, et al. The risk of bias in observational studies of exposures (ROBINS-E) tool: concerns arising from application to observational studies of exposures. Syst Rev. 2018;7(1):242.
\end{references}
\endgroup
